\section{RELATED WORK}

\paragraph{Few-shot and generative NER.}
Few-NERD established a fine-grained few-shot benchmark, and structured
nearest-neighbor and decomposed meta-learning methods illustrate how support
composition affects limited-label learning
\cite{ding2021fewnerd,yang2020structured,ma2022decomposed}. Learned prompt
retrieval makes the same issue explicit for in-context learning
\cite{rubin2022promptretrieval}. Generative IE extends extraction through
structured outputs and instructions
\cite{xu2024llmiesurvey,lu2022uie,wang2023instructuie,wang2023gptner}.
Schema-guided systems such as OneKE, GoLLIE, UniversalNER, and GLiNER further
show the value of explicit task or type descriptions
\cite{luo2025oneke,sainz2024gollie,zhou2024universalner,zaratiana2024gliner}.
EC-FSIE complements these modeling approaches by controlling the provenance
and isolated effects of evidence supplied to a fixed local extractor.

\paragraph{Terminology, retrieval, and verification.}
Gazetteers provide established lexical support for NER
\cite{liu2019gazetteers}; EC-FSIE makes this support query-matched and bounded.
RAG grounds generation in retrieved passages, while retrieval and reranking
choices determine what evidence reaches the model
\cite{lewis2020rag,gao2023ragsurvey}. SelfCheckGPT and Chain-of-Verification
motivate explicit verification, and RAGTruth shows that retrieved context does
not itself eliminate unsupported output
\cite{manakul2023selfcheckgpt,dhuliawala2024cove,niu2024ragtruth}. Our evaluated
NER chain therefore isolates RAG as an evidence channel and interprets its
precision--recall effect directly; verification remains a separate boundary.

\paragraph{Chinese equipment-domain extraction.}
MilitaryCorpus studies fuzzy military-entity boundaries, while work on
equipment-support text, weapon entities, and equipment-fault records develops
task-specific sequence-labeling protocols
\cite{yin2020militaryner,liu2021equipmentsupport,guo2023militaryweapons,gao2024equipmentfaultner}.
These studies establish both the domain need and its annotation difficulty.
EC-FSIE asks a different applied question: how to make demonstration,
terminology, retrieval, and execution controls separately observable in
low-label generative NER.

\paragraph{Resources and positioning.}
CLUENER is a public Chinese NER benchmark, DuIE 2.0 supports Chinese relation
extraction, and CMNEE is a document-level military event-extraction resource
\cite{xu2020cluener,baidu2026duie,zhu2024cmnee}. Their schemas differ from the
four equipment-domain types studied here. We use CLUENER only as a second,
schema-native evaluation of the A/B/C/D control sequence. Thus, the paper's
distinct contribution is not a new universal extractor, but a reproducible
evidence-control framework and an operating-point analysis for low-label
Chinese equipment NER.
